\section{讨论}
\label{sec:discussion_zh}

\subsection{社会状态如何作用于市场}

市场反馈始于持仓和信息曝光的差异。相近的风险敞口提高既有同伴关系的相对权重。消息送达后，其影响取决于局部感染状态，以及主体信心如何调节个人判断与同伴信号的权重。由此形成的委托在订单簿中与其他委托相遇。结算改变下一步可见的持仓和价格参考点。共同父世界分支将干预前的市场历史固定下来，使这条作用路径可以逐步观测。

在该组比较中，传播和从众改变了不同的市场结果。关闭二次感染后，两条价格路径仍整体上行，但回撤幅度下降。关闭从众后，决定同一最大回撤的谷值价格相同，成交量却已出现差异。结合成交与价格记录，就能在价格路径分化前识别交易活动的变化。反转回顾性估值的符号使价格末值低于起点。结合感染、信心和成交记录，可以区分各机制开关产生的响应；这些效应以给定父世界状态和事件为条件。

\subsection{交易制度决定市场参与条件}

RQ1 的双边盘口覆盖率界定了观察价差所需的交易条件。RQ2 的停牌分支把同样的问题带到成交价上。零回撤和第 20 步的价格确认都来自停牌期间沿用的成交价，同一期间，感染与情绪仍在演化。因此，依据价格判断市场是否消化了冲击，需要结合当时的交易机会与社会状态。其余分支展示三个不同的制度环节。价带限定可接受的报价，交收锁仓改变可卖股数，卖方税在结算时重新分配现金。本次记录中，报价处于价带之内，股份锁定未改变实际卖单，税费在成交相同的条件下改变账本余额。因此，比较制度效果时可以结合委托准入、资金和股份预留、现金结算、价格及换手率。

\subsection{LLM 决策取决于观察时序}

系统将语言模型决策接入规则智能体所用的同一个市场。LLM 可以依据可见消息和账户事实生成委托；类型化意图和整数资源检查使委托遵守统一撮合规则。公开发布使公告进入市场视图。要把公告发送给模型，还需要主体处于活跃状态且当前步符合请求计划。离线核对显示，即使公告落在周期请求步，一半种子的慢组依然没有主体同时满足这两个条件。资源自适应缓存可在请求之后复用交易方向与预算比例。评估这一策略需要同时记录发给模型的观察，以及后来由缓存重新计算的委托。

\subsection{证据范围}

实测路径覆盖一个单资产市场、三个匹配的日内种子，以及一个父世界的干预分支。多日收益分布需要在更长轨迹上使用留出的市场数据验证，独立父世界可用于估计分支效应的变异。双速策略的行为保真度与实际计算开销，需要在状态及轨迹字段一致的闭环中比较。本次固定观察与离线调度基准分别给出接口响应和各类调度提供的观察机会。
